\scanpage{198}
\mspage{72}
puisque les $i_{\beta'}$ pour $\beta_0\leq\beta'<\beta$ sont \inserted{tous} dans $J_{\beta_0}\supset J_{\beta'}$, donc aussi leur limite, \deleted{\unsure{donc}} \deleted{$i_\beta$ \unsure{est}} donc
\[
 i_\beta\in\bigcap_{\beta_0<\beta}J_{\beta_0}=J_\beta .
\]
Bon la récurrence transfinie marche, et nous donne donc un \unsure{élmt|élément} \deleted{de} \unsure{$s_\alpha$} de $J_\alpha$ !

Je vais prouver aussi la \unsure{forme|formule} $\unsure{\cup|\mathrm U}$ \unsure{assertion} pour $\bigcup_{\alpha\in A}\deleted{I}M_\alpha$, quand $A$ est fini. Posant cette fois $M=\bigcup_\alpha M_\alpha$, il faut prouver, moyennant les relations $(\ast)$\inserted{\unsure{(p. 68)}}, que
\[
 \begin{aligned}
 x\in M &\Longleftrightarrow
  \text{l'ens }I_M\ (= \bigcup_\alpha I_{M_\alpha})\\
 \substack{\text{(i.e. }\exists\alpha\text{ avec}\\ x\in M_\alpha\text{)}} &\hspace{2em}\text{est cofinal dans }I
 \end{aligned}
 \tag*{$(\ast\ast)^{\sim}$}
\]
et uniquement \unsure{car|en} le premier terme de cette \unsure{équivalence} : il \unsure{doit|devrait} signifier que \deleted{$\unsure{\forall}$} $\exists\alpha$ avec $I_{M_\alpha}$ cofinal dans $I$. Bien sûr, si un des $I_{M_\alpha}$ est cofinal dans $I$, \deleted{\unsure{sa}} \inserted{\unsure{à|a} fortiori} leur réunion. Mais inversement, si la réunion est cofinale, elle est filtrante et réunion des parties (filtrantes) $I_{M_\alpha}$, dont l'une doit \unsure{bien|donc} être cofinale dans elle par le lemme 5.14, donc aussi dans $I$, qed.
% archive sheet number: 197
